\documentclass[11pt]{article}
\usepackage[margin=1in]{geometry}
\usepackage[T1]{fontenc}
\usepackage{lmodern,microtype}
\usepackage{amsmath,amssymb,amsthm,mathtools}
\usepackage{booktabs,array}
\usepackage{xurl}
\usepackage[colorlinks=true,linkcolor=blue,citecolor=blue,urlcolor=blue]{hyperref}
\setlength{\emergencystretch}{2em}
\newtheorem{theorem}{Theorem}[section]
\newtheorem{proposition}[theorem]{Proposition}
\newtheorem{lemma}[theorem]{Lemma}
\newtheorem{corollary}[theorem]{Corollary}
\theoremstyle{definition}
\newtheorem{definition}[theorem]{Definition}
\newtheorem{hypothesis}[theorem]{Hypothesis}
\theoremstyle{remark}
\newtheorem{remark}[theorem]{Remark}
\numberwithin{equation}{section}
\newcommand{\Nm}{\mathrm N}
\newcommand{\OO}{\mathcal O}
\newcommand{\CC}{\mathcal C}
\newcommand{\HH}{\mathsf H}
\newcommand{\ind}{\mathbf1}
\newcommand{\rad}{\operatorname{rad}}
\newcommand{\Rea}{\operatorname{Re}}
\newcommand{\ee}{\varepsilon}
\newcommand{\st}{\boldsymbol t}
\newcommand{\conj}[1]{\overline{#1}}
\newcommand{\Tlarge}{\mathcal T_{\mathrm{large}}}
\title{Conductor localization of a mixed inverse--plain moment\\
in a sextic character family}
\author{Drafted for Edward Baker}
\date{Working draft --- 8 October 2026}

\begin{document}
\maketitle
\begin{abstract}
We study a mixed second moment formed from a M\"obius inverse polynomial,
a plain ideal polynomial, and a product of disjoint physical prime sums
over the Eisenstein integers. Under explicitly stated buffered-bin,
marked-inverse, and row-count inputs, we reduce this moment to one signed
sum. Rows of smaller primitive conductor, the entire plain-variable
diagonal, and column pairs of ratio conductor at most $U^{4/5}$ already
fit below the desired exponent. The key elementary estimate counts
$O_\ee(XV^{1/2+\ee})$ pairs of ideals of norm $O(X)$ whose sextic
ratio character has conductor norm at most $V$. The reduction preserves
all zero-extension masks and the actual M\"obius divisor coefficients.
For the parameter range considered here, the discarded terms have
exponent margin at least $1/6250$. A further reduction keeps the row
selector and removes small conductors of the plain-variable ratio,
while increasing the removable total-conductor range. For the source's
physical prime slots, a positive interpolation also proves the mixed
bound when the plain length exponent is at least $21/50$.
We identify the remaining signed
correlation and explain why a direct use of the existing marked inverse
recursion does not estimate it. A principal-row test exposes an unmet
budget for complete-family enlargement, while an elementary mean-zero
profile subcase controls its sixth-power layer. We also calculate the conditional
row-count consequence of a $1/5000$ saving in this mixed moment.
The remaining signed estimate is open; no new zero-free region is
established.
\end{abstract}

\section{Introduction and status}

The character-family method in the September 30 companion preprint
of OpenAI \cite{openai} produces simultaneous inverse and plain
polynomial witnesses for a buffered zero bin. It bounds the resulting
exceptional rows using separate marked moments. A natural next question
is whether the dependence between the two witnesses supplies an
additional saving in their joint moment.

The present draft carries out the first coefficient-level reduction
of that question. Its purpose is to separate contributions that already
follow from the existing inputs from the signed arithmetic estimate
that remains to be proved. The main result is an identity with a
quantitatively smaller error, rather than a proof of the desired new
moment bound.

There are two distinct conductors. The first is the primitive conductor
of the character indexed by a physical row. The second is the conductor
of the ratio of two column characters after the moment is expanded.
Keeping both of them visible removes more of the moment than a
diagonal calculation alone. The exact ratio identity also retains a
coprimality mask at primes where the character phases cancel.

The analytical statements imported from \cite{openai} are assumptions
of the application in this draft. We do not independently verify that
preprint's reflection, detector, moment, or global zero-free theorems.
Its numbering refers to the September 30 paper entitled
\emph{The Quasi-Riemann Hypothesis: A Zero-Free Half-Plane
$\Rea s>7/8$}, not to the distinct October 5 file
\texttt{paper2.pdf}. The elementary pair-counting argument below is
independent of those deep inputs.

The surrounding investigation records the conditional candidate
$7/8-1/24000$ \cite{records}. A further proposed target,
$7/8-1/20000$, would require the new signed estimate isolated here,
as well as the other geometric and contour inputs of that investigation.
Neither candidate is asserted as an independently established
zero-free theorem in this manuscript.

\section{Arithmetic setting and analytic inputs}\label{sec:setting}

\subsection{Characters, profiles, and the actual product}

Let $F=\mathbb Q(\sqrt{-3})$, let $\OO$ be its ring of integers, and
write $\Nm\mathfrak a$ for the absolute norm of a nonzero integral
ideal. Ideal counting gives
\begin{equation}\label{eq:ideal-count}
\#\{\mathfrak a:\Nm\mathfrak a\le Y\}\ll_F Y\qquad(Y\ge1).
\end{equation}
Fix a finite excluded set $\mathcal S$ containing the primes above $6$
and all primes required by the fixed ray data. We suppress fraktur
letters when no confusion is possible. All column ideals and prime
supports below are outside $\mathcal S$.

For a good prime ideal $p$, let $\chi_p(u)$ be its sextic residue
character on element rows, with exact order six on residue-field units,
extended by zero when $p\mid u$. Extend multiplicatively in the ideal
denominator:
\[
\chi_n(u)=\prod_{p\mid n}\chi_p(u)^{v_p(n)}.
\]
A positive power is still zero on nonunits; it is not replaced there
by a unit-valued inverse. Fix one orientation $s_\chi\in\{1,-1\}$ and
one fixed finite-order ray presentation $\nu$, including its original
zeros. The common row presentation is
\begin{equation}\label{eq:presentation}
\psi_u(n)=\nu(n)\chi_n(u)^{s_\chi}.
\end{equation}
The negative orientation has the same zero convention.

Let $U\ge2$ be the physical row scale. The set $\CC$ consists of
nonprincipal physical element rows with $\Nm(u)\asymp U$ in a fixed
buffered zero/profile bin. In the application these rows are
sixth-power-free and retain their unit factors. It is not necessary
that a prescribed pair of witness lengths be large on every row of
$\CC$.

Fix
\begin{equation}\label{eq:ranges}
\frac9{25}\le\delta\le\frac{21}{50},\qquad
\frac7{10}\le r\le\frac{37}{50},\qquad
\frac9{25}\le m\le\frac12.
\end{equation}
Put $D=U^r$, $N=U^m$. Let $I$ be a fixed finite set of whole physical
prime slots, with pairwise disjoint supports $\mathcal P_i$ independent
of $u$, scale $P_i=U^{w_i}$ with $w_i\ge0$ in fixed bounded ranges, and
coefficients $|a_i(p)|\le1$. The slot count is bounded independently
of $U$.
Write $z=\sum_{i\in I}w_i$. The selected slots satisfy fixed strict
marked-inverse margins
\begin{equation}\label{eq:widths}
r+2z\le1-c_1,\qquad 2r+8z\le3-c_2,\qquad c_1,c_2>0.
\end{equation}

All smooth profiles are supported in fixed compact subintervals of
$(0,\infty)$ and have the permitted uniform seminorm bounds. At fixed
separating parameters define
\begin{align}
M_r(u)&=D^{-1/2}\sum_d\mu_F(d)\psi_u(d)A(\Nm d/D),\label{eq:inverse}\\
S_m(u)&=N^{-1/2}\sum_k\psi_u(k)B(\Nm k/N),\label{eq:plain}\\
Q_I(u)&=\prod_{i\in I}\left( P_i^{-1/2}
\sum_{p\in\mathcal P_i}a_i(p)\psi_u(p)W_i(\Nm p/P_i)\right).
\label{eq:slots}
\end{align}
The product in \eqref{eq:slots} means the product of the displayed
slot sums. The empty product equals one. The profile $A$ retains the
actual truncated-inverse cutoff, when present. The common inverse/plain
norm twist is included in $A,B$; physical prime twists remain in $W_i$.
These are the actual profiles, not arbitrary arithmetic coefficients.

Write $\HH=1+T_1$ for the source's height allowance. An exponent denoted
by $A_0$ below is finite and may increase from line to line. Implied
constants may depend on the fixed arithmetic data, supports, slot
count, seminorm orders, strict margins, and requested positive loss.
No uniformity over all target conductors is asserted.

\subsection{Three explicit inputs}

Let $Q_{\psi_u}$ be the norm of the conductor of the primitive
character inducing \eqref{eq:presentation}, and put
\[
\theta(u)=\frac{\log Q_{\psi_u}}{\log U}.
\]
The source gives $Q_{\psi_u}\ll U$; thus
$\theta\le1+O(1/\log U)$, rather than necessarily $\theta\le1$.
We use the following inputs, uniformly in the stated parameter ranges
and permitted profile families.

\begin{hypothesis}[Marked inverse input]\label{hyp:inverse}
For every $\ee>0$, the original row family satisfies
\begin{equation}\label{eq:inverse-input}
\sum_{0<\Nm(u)\le C_0U}|M_r(u)Q_I(u)|^2
\ll_\ee U^{1+\ee}\HH^{A_0},
\end{equation}
where $C_0$ is a fixed bound containing the physical dyad and
\eqref{eq:widths} holds.
\end{hypothesis}

This is the application of \cite[Lemma 17.1]{openai}. Its common
presentation and orientation, disjoint prime supports, and coefficient
requirements are part of our setting. A bounded row restriction may be
dropped after a nonnegative norm has been formed.

\begin{hypothesis}[Buffered plain input]\label{hyp:plain}
For every $\ee>0$ and every $u\in\CC$,
\begin{equation}\label{eq:plain-input}
|S_m(u)|^2\ll_\ee
U^{\delta\min\{m,\theta(u)-m\}+\ee}\HH^{A_0}.
\end{equation}
\end{hypothesis}

Section~\ref{sec:row-conductor} explains how this follows from the
source's buffered reflection argument with the primitive conductor
retained.

\begin{hypothesis}[Full-bin count]\label{hyp:count}
For every $\ee>0$,
\begin{equation}\label{eq:count-input}
\#\CC\ll_\ee U^{9/8-23\delta/20+\ee}\HH^{A_0}.
\end{equation}
\end{hypothesis}

This is a count of the full fixed bin, before selection of the
particular $r,m$ in \eqref{eq:ranges}. In the application it is a
consequence of the previously available scalar row count, as shown
in Section~\ref{sec:application}. It does not assume the new mixed
moment.

The main theorem uses only these three inputs and the fixed coefficient
structure above. The more restrictive amplitude-profile condition
used to locate the last difficult bins is not a hypothesis of the
arithmetic reduction.

\section{Main reduction}\label{sec:main}

Put
\[
\eta=\frac1{5000},\qquad
K=1+\delta m-\eta,\qquad
X=DN\prod_{i\in I}P_i.
\]
For a tuple $\st=(d,k,(p_i)_{i\in I})$ on the original supports define
\begin{align}
n(\st)&=dk\prod_{i\in I}p_i,\label{eq:column}\\
b(\st)&=\mu_F(d)\nu(n(\st))A(\Nm d/D)B(\Nm k/N)
\prod_{i\in I}a_i(p_i)W_i(\Nm p_i/P_i).\label{eq:tuple-coeff}
\end{align}
In particular,
\begin{equation}\label{eq:product}
M_r(u)S_m(u)Q_I(u)
=X^{-1/2}\sum_{\st}b(\st)\chi_{n(\st)}(u)^{s_\chi}.
\end{equation}

For two ideals $n,n'$ set
\begin{align}
e_p(n,n')&=(v_p(n)-v_p(n'))\bmod6\in\{0,\ldots,5\},\nonumber\\
c(n,n')&=\prod_{e_p\ne0}p^{e_p},\qquad
f(n,n')=\rad(c(n,n')),\label{eq:ratio-data}\\
E(n,n')&=\prod_{\substack{p\mid nn'\\e_p=0}}p.\nonumber
\end{align}
The moving conductor is $f$; $E$ records the remaining zeros.
Let
\begin{equation}\label{eq:high-rows}
\CC_+=\{u\in\CC:Q_{\psi_u}>U^{2m-1/1000}\},
\end{equation}
and define the restricted kernel
\begin{equation}\label{eq:kernel}
\mathcal K_{\CC_+}(c;E)=
\sum_{u\in\CC_+}\chi_c(u)^{s_\chi}\ind_{(u,E)=1}.
\end{equation}
In the following sum, $n=n(\st)$ and $n'=n(\st')$:
\begin{equation}\label{eq:Tlarge}
\Tlarge=\frac1X
\sum_{\substack{\st,\st'\\k\ne k'\\
\Nm f(n,n')>U^{4/5}}}
b(\st)\conj{b(\st')}\,
\mathcal K_{\CC_+}\bigl(c(n,n');E(n,n')\bigr).
\end{equation}

\begin{theorem}[Conductor-localized mixed moment]\label{thm:main}
Under the setting of Section~\ref{sec:setting} and
Hypotheses~\ref{hyp:inverse}--\ref{hyp:count}, for every $\ee>0$,
\begin{equation}\label{eq:main}
\sum_{u\in\CC}|M_r(u)S_m(u)Q_I(u)|^2
=\Rea\Tlarge+
O_\ee\!\left(U^{K-1/6250+\ee}\HH^{A_0}\right).
\end{equation}
The identity and estimate first hold at fixed separating parameters.
Their application to rowwise choices requires the same uniform
estimates for the derivative profiles in the source's Sobolev procedure.
\end{theorem}

The theorem removes three disjoint contributions. Their exponent
margins below $K$ are as follows:
\begin{center}
\begin{tabular}{@{}p{0.69\textwidth}c@{}}
\toprule
Contribution & Margin\\
\midrule
Rows with $Q_{\psi_u}\le U^{2m-1/1000}$ & $1/6250$\\
Plain-variable diagonal $k=k'$ on the remaining rows & $647/5000$\\
Remaining pairs with $\Nm f(n,n')\le U^{4/5}$ & $23/1250$\\
\bottomrule
\end{tabular}
\end{center}

\begin{corollary}[The exact open input]\label{cor:open}
To prove
\begin{equation}\label{eq:desired}
\sum_{u\in\CC}|M_rS_mQ_I|^2
\ll_\ee U^{1+\delta m-1/5000+\ee}\HH^{A_0},
\end{equation}
it suffices to prove, with the required uniformity,
\begin{equation}\label{eq:open}
\Rea\Tlarge\le C_\ee U^{1+\delta m-1/5000+\ee}\HH^{A_0}.
\end{equation}
Conversely, \eqref{eq:desired} and the theorem imply
\eqref{eq:open}, after adjusting the harmless constants and losses.
\end{corollary}

The remaining block need not be positive. Theorem~\ref{thm:main}
also bounds its negative part by the smaller error because the complete
moment is nonnegative. Since the complete pair sum is real, a modulus
bound for the whole sum \eqref{eq:Tlarge} is therefore equivalent at
the target scale and is also sufficient. Summing separate absolute
bounds for its individual pairs asks for a stronger estimate; no such
replacement is made here.

\section{Exact kernels and an elementary pair count}\label{sec:pairs}

\begin{lemma}[Ratio character with its zeros]\label{lem:ratio}
For every element row $u$,
\begin{equation}\label{eq:ratio-identity}
\chi_n(u)^{s_\chi}\conj{\chi_{n'}(u)^{s_\chi}}
=\chi_{c(n,n')}(u)^{s_\chi}\ind_{(u,E(n,n'))=1}.
\end{equation}
At good primes the character on the right has moving primitive
conductor exactly $f(n,n')$.
\end{lemma}
\begin{proof}
On units at $p$, reducing the difference of valuations modulo six
gives the same local phase. If that residue is nonzero, a nontrivial
power of the exact-order-six residue character is nonprincipal modulo
$p$, so its conductor is $p$. If the residue is zero, its phase is
trivial, but the original product still vanishes when $p\mid u$
whenever $p$ divides either original column. Those are exactly the
primes in $E$. The Chinese remainder theorem gives the product
identity and conductor assertion. Any fixed ray factors introduced
by a presentation are kept separately; they do not cancel ramification
at a good prime.
\end{proof}

For example, $n=p^6$ and $n'=1$ give the factor $\ind_{p\nmid u}$,
not the constant one. Removing such a mask would change the moment.

\begin{lemma}[Coefficient majorant]\label{lem:coeff}
Let
\[
B_*(n)=\sum_{\st:n(\st)=n}|b(\st)|.
\]
The support has $cX\le\Nm n\le CX$ for fixed positive constants,
and, for every $\ee>0$,
\begin{equation}\label{eq:coeff-bound}
B_*(n)\ll_\ee(\Nm n)^\ee,\qquad
X^{-1}\sum_n B_*(n)^2\ll_\ee X^\ee.
\end{equation}
The implied constants have the allowed profile dependence.
\end{lemma}
\begin{proof}
All normalized norms lie in fixed annuli. Each factor in
\eqref{eq:tuple-coeff} has bounded absolute value there. A fixed number
of ordered ideal factors has divisor-bounded multiplicity; imposing
prime supports, squarefreeness of $d$, or other original zeros only
reduces that multiplicity. Thus $B_*(n)$ is bounded by a fixed power
of the ideal divisor function, which is $O_\ee((\Nm n)^\ee)$.
Apply \eqref{eq:ideal-count}, dividing the preliminary loss as needed,
to obtain the second assertion.
\end{proof}

\begin{lemma}[Pairs of small ratio conductor]\label{lem:pair-count}
For fixed $C>0$, $X\ge1$, $V\ge1$, and every $\ee>0$,
\begin{equation}\label{eq:pair-count}
\#\{(n,n'):\Nm n,\Nm n'\le CX,\ \Nm f(n,n')\le V\}
\ll_{C,\ee} X V^{1/2+\ee}.
\end{equation}
\end{lemma}
\begin{proof}
Write $g=(n,n')$, $n=gA$, $n'=gB$, with $(A,B)=1$. Extract sixth
powers uniquely:
\[
A=a_0a^6,\qquad B=b_0b^6,
\]
where $a_0,b_0$ are sixth-power-free and coprime. Then
$\rad(a_0b_0)=f(n,n')$. At each prime of $f$ there are ten allocations:
one of two sides and one of the exponents $1,\ldots,5$.

Fix an allocation and the ideals $a,b$. Ideal counting bounds the
number of possible $g$ by
\begin{align}
&\ll \frac{X}{\max\{\Nm a_0(\Nm a)^6,\Nm b_0(\Nm b)^6\}}\nonumber\\
&\le \frac{X}{\sqrt{\Nm(a_0b_0)}\,(\Nm(ab))^3}.
\label{eq:g-bound}
\end{align}
When the allowed norm for $g$ is below one, the actual count is zero
and this positive majorant is still valid. Dropping restrictions on
$a,b$ enlarges the bound. Their ideal sums converge at exponent three.
Since $\Nm(a_0b_0)\ge\Nm f$, the total for fixed $f$ is
\[
\ll X\,10^{\omega(f)}(\Nm f)^{-1/2}.
\]
The divisor bound $10^{\omega(f)}\ll_\ee(\Nm f)^\ee$, followed by
partial summation in \eqref{eq:ideal-count}, gives
\[
\sum_{\substack{f\ {\rm squarefree}\\\Nm f\le V}}
10^{\omega(f)}(\Nm f)^{-1/2}\ll_\ee V^{1/2+\ee}.
\]
This proves the claim.
\end{proof}

\begin{proposition}[Absolute small-conductor block]\label{prop:small}
Let $\mathcal D$ be any subset of the physical rows. In the expansion
of the mixed norm over $\mathcal D$, the sum of absolute tuple weights
with $\Nm f(n,n')\le V$ is at most
\begin{equation}\label{eq:small-block}
\ll_\ee \#\mathcal D\,V^{1/2+\ee}X^\ee.
\end{equation}
The same bound holds after any additional tuple restriction, including
$k\ne k'$.
\end{proposition}
\begin{proof}
The restricted row kernel has absolute value at most $\#\mathcal D$.
Group the positive tuple majorant by $n,n'$, apply
Lemma~\ref{lem:coeff}, then Lemma~\ref{lem:pair-count}. An additional
tuple restriction can only reduce this positive majorant. It is not
necessary for the restricted coefficient to factor as a product of
two composite-column coefficients.
\end{proof}

In particular, $V=1$ controls all sixth-power-equivalent column
pairs, including the ordinary diagonal, by
$O_\ee(\#\mathcal D\,X^\ee)$. The proof does not use cancellation in
the row kernel and is valid for a zero/profile-selected set.

\section{The three discarded contributions}\label{sec:discard}

\subsection{Primitive row conductor}\label{sec:row-conductor}

We first record the source justification of Hypothesis~\ref{hyp:plain}.
Let $a=(1+\delta)/2$ be the buffered zero-bin parameter. In the proof
of \cite[Lemma 8.1]{openai}, the primitive functional equation on
the reflected line gives, before replacing its conductor by $U$,
\begin{equation}\label{eq:reflected-L}
|L_{\rm orig}(1-a-6e+it,\psi_u)|
\ll Q_{\psi_u}^{\,a-1/2+6e}
U^{O(e)+\ee}\HH^{A_0}.
\end{equation}
The conjugate primitive $L$-value is controlled in the buffered
rectangle. The original deleted Euler product is retained. Its radical
has norm $O(U)$, and here $1-a-6e\ge29/100-6e>0$ for sufficiently
small fixed $e$. Its absolute product has an arbitrarily small
$U$-power on this line. The weaker $U^{O(e)+\ee}$ notation in
\eqref{eq:reflected-L} is sufficient.

Repeating the finite-height Mellin shift of
\cite[Lemma 8.2]{openai} yields the reflected bound
$U^{\delta(\theta-m)+\ee}\HH^{A_0}$ for $|S_m|^2$.
The right-hand shift gives $U^{\delta m+\ee}$.
The same cumulative height allowance, buffered joins, and fixed
tail orders are required. The tail order can be chosen to dominate
even a negative displayed exponent. These two bounds give
\eqref{eq:plain-input}.

Let $\CC_-=\CC\setminus\CC_+$. On $\CC_-$, the reflected estimate is
\[
|S_m|^2\ll_\ee U^{\delta m-\delta/1000+\ee}\HH^{A_0}.
\]
After this pointwise estimate, positivity permits enlargement of the
remaining marked inverse norm to the full row family in
Hypothesis~\ref{hyp:inverse}. Therefore
\begin{equation}\label{eq:low-row}
\sum_{u\in\CC_-}|M_rS_mQ_I|^2
\ll_\ee U^{K-1/6250+\ee}\HH^{A_0},
\end{equation}
because
\[
\frac{\delta}{1000}-\frac1{5000}
\ge\frac{9}{25000}-\frac1{5000}=\frac1{6250}.
\]
This does not replace the physical row scale by $U^\theta$.

\begin{remark}[Actual detector spikes]
If a selected plain witness also satisfies
$|S_m|^2\gg U^{\delta m-\ee_w}$, then the reflected bound forces
\[
\theta\ge2m-\frac{\ee+\ee_w}{\delta}
-\frac{A_0\log\HH}{\delta\log U}-o(1).
\]
After the permitted loss and height choices this reads
$\theta\ge2m-o_{\rm loss}(1)$. The loss notation describes adjustable
fixed choices, not a new uniform zero-loss assertion. For the larger
moment set $\CC$, the positive bound \eqref{eq:low-row} is the
appropriate argument; the witness lower bound need not hold on all rows.
\end{remark}

\subsection{The complete plain-variable diagonal}

Expand only the plain factor over $\CC_+$. The terms with $k=k'$
form the exact nonnegative expression
\begin{equation}\label{eq:plain-diagonal}
D_{\rm plain}=
\frac1N\sum_k |B(\Nm k/N)|^2|\nu(k)|^2
\sum_{u\in\CC_+}|M_r(u)Q_I(u)|^2\ind_{(u,k)=1}.
\end{equation}
The fixed character zeros are written explicitly as $|\nu(k)|^2$.
They equal one on the original good support; retaining them also
handles any additional fixed zero. After forming the positive sum
we may drop its coprimality restriction. The normalized coefficient
square mass is $O(1)$ by annular ideal counting. Hence
\begin{equation}\label{eq:diag-bound}
0\le D_{\rm plain}\ll_\ee U^{1+\ee}\HH^{A_0}
\ll_\ee U^{K-647/5000+\ee}\HH^{A_0},
\end{equation}
since $\delta m-1/5000\ge647/5000$.
This removes all inverse and prime cross terms paired with equal
plain ideals, not merely equal composite columns.

\subsection{Small ratio conductors and proof of the theorem}

Among the remaining tuples $k\ne k'$, let $B_{\rm small}$ denote
the signed contribution with $\Nm f(n,n')\le U^{4/5}$.
Proposition~\ref{prop:small}, Hypothesis~\ref{hyp:count}, and the bounded
range $X=U^{r+m+z}$ imply
\begin{align}
|B_{\rm small}|
&\ll_\ee U^{9/8-23\delta/20+2/5+\ee}\HH^{A_0}\nonumber\\
&\ll_\ee U^{K-23/1250+\ee}\HH^{A_0}.
\label{eq:small-margin}
\end{align}
The last inequality follows from
\begin{equation}\label{eq:margin-arithmetic}
\delta\left(m+\frac{23}{20}\right)-\frac{21}{40}-\frac1{5000}
\ge\frac9{25}\left(\frac9{25}+\frac{23}{20}\right)
-\frac{21}{40}-\frac1{5000}=\frac{23}{1250}.
\end{equation}
The small divisor losses can be divided among the preceding estimates
because all length exponents remain in fixed bounded intervals.

\begin{proof}[Proof of Theorem~\ref{thm:main}]
Partition the rows into $\CC_-$ and $\CC_+$. On $\CC_+$ expand the
plain factors and remove $k=k'$ as in \eqref{eq:plain-diagonal}.
For the remaining tuple pairs, use Lemma~\ref{lem:ratio} and partition
by $\Nm f\le U^{4/5}$ or $\Nm f>U^{4/5}$. This gives exactly
\[
\sum_{u\in\CC}|M_rS_mQ_I|^2
=\sum_{u\in\CC_-}|M_rS_mQ_I|^2
+D_{\rm plain}+B_{\rm small}+\Tlarge.
\]
The pair regions are invariant under interchanging the two tuples,
so their complete sums are real. The first three terms have the
bounds \eqref{eq:low-row}, \eqref{eq:diag-bound}, and
\eqref{eq:small-margin}. Their smallest guaranteed margin is
\[
\min\left\{\frac1{6250},\frac{647}{5000},\frac{23}{1250}\right\}
=\frac1{6250}.
\]
Taking absolute values only for these discarded terms proves
\eqref{eq:main}. A finite enlargement of $A_0$ accommodates their
height factors.
\end{proof}

\section{A refinement with two ratio conductors}\label{sec:two-ratios}

Apply the definitions \eqref{eq:ratio-data} to the plain ideals $k,k'$
alone and denote the resulting data by
$c_{\rm pl},f_{\rm pl},E_{\rm pl}$. These are distinct from the
total-column data $c(n,n'),f(n,n'),E(n,n')$ and from the primitive
conductor of a physical row. The inverse and prime variables can
introduce or cancel local phases, so neither ratio conductor generally
determines the other.

\begin{lemma}[Weighted plain pairs]\label{lem:weighted-plain}
Let $w(u)\ge0$ have finite mass $\mathcal A=\sum_u w(u)$.
Write $a(k)=\nu(k)B(\Nm k/N)$ with its original zeros, and set
\[
\mathcal K_w^{\rm pl}(k,k')=
 \sum_u w(u)\chi_{c_{\rm pl}(k,k')}(u)^{s_\chi}
 \ind_{(u,E_{\rm pl}(k,k'))=1}.
\]
For $V\ge1$,
\begin{equation}\label{eq:weighted-plain}
\frac1N\sum_{\Nm f_{\rm pl}(k,k')\le V}
 |a(k)a(k')\mathcal K_w^{\rm pl}(k,k')|
\ll_\ee \mathcal A V^{1/2+\ee}N^\ee\HH^{A_0}.
\end{equation}
The same bound holds with additional plain-pair restrictions.
\end{lemma}
\begin{proof}
The exact local identity in Lemma~\ref{lem:ratio} retains the displayed
mask. Each zero-extended character has modulus at most one, so the
weighted kernel is bounded by $\mathcal A$. Apply
Lemma~\ref{lem:pair-count} at column scale $N$, and use the bounded
annular coefficients. Additional restrictions reduce this positive
majorant.
\end{proof}

Take $w(u)=\ind_{\CC_+}(u)|M_r(u)Q_I(u)|^2$. Its mass is
$O_\ee(U^{1+\ee}\HH^{A_0})$ by Hypothesis~\ref{hyp:inverse}.
The entire signed low-plain-conductor contribution, including the
plain diagonal, is therefore $O_\ee(U^{1+\ee}V^{1/2+\ee}\HH^{A_0})$.
This estimate assembles the inverse and prime variables into a positive
weight before taking absolute values; it is not a termwise bound for
their expanded tuples.

Put $\gamma=1/6250$, and define
\begin{align}
v_{\rm pl}&=2\delta m-\frac9{12500},\qquad
 V_{\rm pl}=U^{v_{\rm pl}},\label{eq:plain-threshold}\\
\tau_0&=2\delta\left(m+\frac{23}{20}\right)
 -\frac14-\frac9{12500},\qquad
 V_{\rm tot}=U^{\tau_0}.\label{eq:total-threshold}
\end{align}
Since $2(\eta+\gamma)=9/12500$, these satisfy
\[
1+v_{\rm pl}/2=K-\gamma,\qquad
\frac98-\frac{23\delta}{20}+\tau_0/2=K-\gamma.
\]
Their ranges on \eqref{eq:ranges} are
\begin{equation}\label{eq:threshold-ranges}
\frac{3231}{12500}\le v_{\rm pl}\le\frac{5241}{12500},\qquad
\frac{2614}{3125}\le\tau_0\le\frac{14191}{12500}.
\end{equation}
In particular $\tau_0\ge0.83648>4/5$. A ratio conductor of columns
is not bounded by the physical row scale, so $\tau_0>1$ causes no
inconsistency.

Define the further restricted, exact signed sum
\begin{equation}\label{eq:Tjoint}
\mathcal T_{\rm joint}=\frac1X
\sum_{\substack{\st,\st'\\
 \Nm f_{\rm pl}(k,k')>V_{\rm pl}\\
 \Nm f(n(\st),n(\st'))>V_{\rm tot}}}
b(\st)\conj{b(\st')}\,
\mathcal K_{\CC_+}\bigl(c(n,n');E(n,n')\bigr).
\end{equation}
The condition $k\ne k'$ is automatic here.

\begin{theorem}[Two-ratio localization]\label{thm:two-ratios}
Under the same hypotheses as Theorem~\ref{thm:main},
\begin{equation}\label{eq:two-ratio-reduction}
\sum_{u\in\CC}|M_rS_mQ_I|^2
=\Rea\mathcal T_{\rm joint}
+O_\ee(U^{K-1/6250+\ee}\HH^{A_0}).
\end{equation}
\end{theorem}
\begin{proof}
Remove $\CC_-$ using \eqref{eq:low-row}. On $\CC_+$ expand only
the plain factor and remove its entire block $\Nm f_{\rm pl}\le
V_{\rm pl}$ using \eqref{eq:weighted-plain}. Expand the remaining
block in the inverse and prime variables. Proposition~\ref{prop:small}
and Hypothesis~\ref{hyp:count} bound the part with
$\Nm f(n,n')\le V_{\rm tot}$ by $O_\ee(U^{K-\gamma+\ee}\HH^{A_0})$;
the extra plain-pair restriction only reduces its positive majorant.
This leaves \eqref{eq:Tjoint}. All pair regions are swap-invariant,
so their complete signed sums are real.
\end{proof}

No total-conductor condition was inserted into the positive inverse
weight. The proof removes disjoint contributions in the stated order.
Small plain-ratio conductor can occur for distinct $k,k'$ while the
inverse variables make the total ratio conductor large; this removes
additional terms from the earlier retained sum. The strengthened open
input is now
\begin{equation}\label{eq:joint-open}
\Rea\mathcal T_{\rm joint}
\le C_\ee U^{1+\delta m-1/5000+\ee}\HH^{A_0}.
\end{equation}
It remains equivalent at this scale to the mixed bound. Positivity of
\eqref{eq:two-ratio-reduction}'s left side bounds the negative part by
the smaller error, so a modulus bound for the whole remaining sum is
also sufficient. No individual-pair cancellation is asserted.

\section{Uniformity and the row selector}\label{sec:uniformity}

The exact coefficient expansion is made at one fixed set of separating
parameters. It is not an identity with a common coefficient array if
each row is assigned an unrelated height or profile inside the
coefficient. The source handles rowwise choices by a fixed Sobolev
argument on a positive norm, using finitely many derivative profiles
and parameter integrals. These derivatives preserve the annular
coefficient structure and add permitted logarithmic or height factors.
Any proof of \eqref{eq:open} intended for that application must be
uniform for the resulting derivative profiles.

The order of choices matters. The requested real exponent losses and
strict capacities are fixed first; the source's internal seminorm and
height orders are then determined; only afterward is the height range
chosen sufficiently slowly growing to absorb its fixed polynomial cost.
The numerical margins above are positive fixed numbers and allow this
ordering. They do not grant arbitrary height uniformity.

The arithmetic indicator of $\CC_+$ depends on zero bins, physical
prime amplitudes, and primitive conductors. It is not a smooth complete
row weight. Consequently a complete-row character-sum estimate does
not become an estimate for $\mathcal K_{\CC_+}$ simply because its
ratio character is nonprincipal. Positivity permits enlarging the
\emph{whole} mixed norm to a complete smooth family. It does not permit
termwise enlargement of a signed off-diagonal block.

A fixed finite sum of presentations or whole products can be handled
by Cauchy--Schwarz before applying the single-presentation estimate.
The exact equality \eqref{eq:main} is stated for the single product
\eqref{eq:product}; it does not omit cross-presentation terms from an
unexpanded finite sum.

\subsection{The principal-row feasibility test}\label{sec:selector-test}

Let $\Phi$ be a nonnegative smooth compactly supported annular
majorant for the physical norm interval. The complete-family route
would require the stronger positive estimate
\begin{equation}\label{eq:complete-enlargement}
\sum_{u\ne0}\Phi(\Nm u/U)|M_r(u)S_m(u)Q_I(u)|^2
\ll_\ee U^{K+\ee}\HH^{A_0}.
\end{equation}
This adds principal and non-sixth-power-free rows. We test the
principal presentation $\nu=1$, retaining its fixed exclusions.
For each integral ideal $a$, choose a generator $\alpha_a$. Since
the Eisenstein ring has six units, $\alpha_a^6$ is independent of
that choice, and different ideals give different sixth powers.
Thus the exact positive sixth-power contribution is
\begin{equation}\label{eq:sixth-layer}
E_6=\sum_a\Phi((\Nm a)^6/U)
 |M_r^{(a)}|^2|S_m^{(a)}Q_I^{(a)}|^2.
\end{equation}
There is no extra factor six. The superscript means that every
original coefficient and fixed exclusion is retained and that the
additional masks are $(a,d)=1$, $(a,k)=1$, and $p_i\nmid a$.
This follows from $\chi_n(\alpha_a^6)=\ind_{(a,n)=1}$.
These rows have primitive inducing conductor one; the mask radical
is not a primitive conductor, nor is it a column-ratio conductor.

Put $\beta_B=\int_0^\infty B(y)\,dy$, including the actual norm
twist. With $R_a$ the radical containing $a$ and the fixed excluded
primes, smooth lattice counting and inclusion--exclusion give
\begin{align}
S_m^{(a)}&=N^{1/2}c(a)\beta_B
 +O_\ee(N^{-1/2}U^\ee\HH^{A_0}),\label{eq:principal-plain}\\
c(a)&=\mathop{\rm Res}_{s=1}\zeta_F(s)
 \prod_{p\mid R_a}(1-(\Nm p)^{-1}).\nonumber
\end{align}
Indeed the unmasked smooth ideal sum has its lattice area term and
an $O(\HH^{A_0})$ error on scales at least one. For $e\mid R_a$,
the residual scale $N/\Nm e$ exceeds one because
$\Nm e\ll U^{1/6}$ and $m\ge9/25>1/6$.
There are $O_\ee(U^\ee)$ such divisors. This also follows from
\cite[Lemma 18.3]{openai}. Each prime factor satisfies the exact formula
\begin{equation}\label{eq:principal-prime}
Q_i^{(a)}=Q_i^{(1)}-P_i^{-1/2}
 \sum_{\substack{p\in\mathcal P_i\\p\mid a}}
 a_i(p)W_i(\Nm p/P_i).
\end{equation}
The deleted sum costs at most $O(P_i^{-1/2}\log U)$ times the
profile bound. Neither $\beta_B$ nor $Q_i^{(1)}$ is assumed nonzero.
Their actual coefficients must be evaluated before assigning a
principal main-term exponent.

With $G(a)=U^{-m-z}|S_m^{(a)}Q_I^{(a)}|^2$, bounding
\eqref{eq:sixth-layer} at the desired scale is exactly the demand
\begin{equation}\label{eq:principal-inverse-budget}
\sum_a\Phi((\Nm a)^6/U)G(a)|M_r^{(a)}|^2
\ll_\ee U^{1-(1-\delta)m-z-\eta+\ee}\HH^{A_0}.
\end{equation}
If the actual plain and prime factors have nonvanishing principal
mass on a specified root subfamily, then $G(a)\ge U^{-\ee}$ there,
after allocating the small losses. The corresponding unweighted
inverse average must then meet this exponent as well.

Assuming the source's seven-eighths theorem, its
Lemmas~4.9--4.10 and smooth Mellin inversion give only the available
pointwise bound $|M_r^{(a)}|^2\ll_\ee D^{3/4+\ee}U^\ee\HH^{A_0}$.
The masked Dirichlet series is
$\zeta_F(s)^{-1}\prod_{p\mid R_a}(1-(\Nm p)^{-s})^{-1}$;
the contour is placed at $\Rea s=7/8+v$, with $v>0$ arbitrarily
small and fixed. Counting the $O(U^{1/6})$ roots gives inverse
average exponent $1/6+3r/4$. At $z=(1-r)/2-\rho$, its excess over
the unweighted budget is
\begin{equation}\label{eq:sixth-budget-gap}
\Delta_6=(1-\delta)m+\frac r4-\frac13+\eta-\rho
\in\left[\frac{19}{375}-\rho,\frac{1289}{7500}-\rho\right].
\end{equation}
This is a deficit in an available upper bound, not a lower bound
for $E_6$ or a contradiction to \eqref{eq:complete-enlargement}.
For general profiles the exact condition remains
\eqref{eq:principal-inverse-budget}.

\begin{proposition}[A mean-zero principal-layer subcase]\label{prop:mean-zero}
If $\nu=1$ and the actual smooth profile has $\beta_B=0$, then,
with all original masks and prime supports retained,
\begin{equation}\label{eq:mean-zero-layer}
E_6\ll_\ee U^{1/6+r-m+z+\ee}\HH^{A_0}
\ll_\ee U^{K-6791/15000+\ee}\HH^{A_0}.
\end{equation}
\end{proposition}
\begin{proof}
Equation~\eqref{eq:principal-plain} gives
$|S_m^{(a)}|^2\ll_\ee N^{-1}U^\ee\HH^{A_0}$.
Absolute ideal counting gives $|M_r^{(a)}|^2\ll D\HH^{A_0}$
and $|Q_I^{(a)}|^2\ll U^z\HH^{A_0}$. Sum over the roots.
Using $z\le(1-r)/2$, the margin below $K$ is at least
$1/3-r/2+(1+\delta)m-\eta\ge6791/15000$.
\end{proof}

This proposition controls one added layer only. Mean zero at one
height does not imply mean zero for its derivative profiles, so it
does not supply the application uniformity by itself. More generally,
principal-inducing rows have finitely many sixth-power-free seeds
supported on the fixed excluded primes, with their unit factors;
the same calculation applies to each seed with its masks retained.
The detailed selector note in the investigation records this
decomposition and the profile-sensitive budget. Until those added
rows and the remaining complete family are controlled, the main
route must retain the selector in \eqref{eq:high-rows}. No new mixed
saving follows from \eqref{eq:mean-zero-layer}.

\section{A selector-preserving mixed subcase}\label{sec:mixed-subcase}

The following partial result uses additional source inputs, beyond the
three hypotheses of Theorem~\ref{thm:two-ratios}. The slots must be the
source's physical prime sums, with their fixed ray-class restrictions,
smooth annuli and fixed finite-ray coefficient expansions. Arbitrary
prime subsets with bounded coefficients do not automatically satisfy
the buffered pointwise prime estimate used here.
Assume the source's seven-eighths theorem so that $\kappa=3/4$ is
legal in \cite[Lemma 18.1]{openai}. Its slot mesh and all original
height conventions are required.

For a subset of whole slots $J\subset I$ with
$2m+(9/2)z_J\le1$, where $z_J=\sum_{i\in J}w_i$, the source gives
\begin{equation}\label{eq:fourth-subset}
\sum_{u\in\CC_+}|S_m(u)^2Q_J(u)|^2
\ll_\ee U^{1+\ee}\HH^{A_0}.
\end{equation}
For $J=\varnothing$, its zero-slot assertion applies without this
length restriction. Every primitive character in the source's fixed
finite exceptional family has bounded conductor. Thus $\CC_+$ is
in the nonexceptional family for all sufficiently large $U$.
The source's Lemmas~8.2 and 19.1 give, for the actual profiles,
\begin{equation}\label{eq:pointwise-subset}
|M_r|^2\ll_\ee U^{\delta r+\ee}\HH^{A_0},\qquad
|Q_{I\setminus J}|^2\ll_\ee
 U^{\delta(z-z_J)+\ee}\HH^{A_0}.
\end{equation}
Alternatively, \eqref{eq:fourth-subset}--\eqref{eq:pointwise-subset}
may be stipulated explicitly as the extra hypotheses for the next
proposition.

\begin{proposition}[Upper plain-length band]\label{prop:upper-band}
Under these additional inputs, suppose the source mesh permits a
whole-slot selection with total capacity decrement and rounding loss
at most $1/1000$. Then, for $21/50\le m\le1/2$,
\begin{equation}\label{eq:upper-band}
\sum_{u\in\CC_+}|M_rS_mQ_I|^2
\ll_\ee U^{K-59/100000+\ee}\HH^{A_0}.
\end{equation}
Consequently the full moment over $\CC$ has the bound
$O_\ee(U^{K-1/6250+\ee}\HH^{A_0})$ in this band.
\end{proposition}
\begin{proof}
The summand factors exactly as
\[
|M_rS_mQ_I|^2
=|M_rQ_{I\setminus J}|\,|M_rQ_I|\,|S_m^2Q_J|.
\]
Bound the first factor pointwise and apply Cauchy to the other two,
using Hypothesis~\ref{hyp:inverse} and \eqref{eq:fourth-subset}.
This proves the fixed-subset bound
\begin{equation}\label{eq:subset-interpolation}
\sum_{\CC_+}|M_rS_mQ_I|^2
\ll_\ee U^{1+\delta(r+z-z_J)/2+\ee}\HH^{A_0}.
\end{equation}
Put $c_m=2(1-2m)/9$. Greedy selection of whole slots inside $I$
gives $z_J\le c_m$ and
$z-z_J\le(z-c_m)_++\ell_{\rm sel}$, where
$\ell_{\rm sel}\le1/1000$ includes the decrement and last slot.
This includes the empty selection when $c_m=0$.
Since $z\le(1-r)/2$ and $(1-r)/2>c_m$ on the box, the saving in
\eqref{eq:subset-interpolation} below $1+\delta m$ is at least
\[
\delta\left(\frac79m-\frac r4-\frac5{36}
 -\frac{\ell_{\rm sel}}2\right).
\]
Without the selection loss it is at least $1/1000$ in the stated
band. The loss is at most $21/100000$. Subtracting $\eta=1/5000$
leaves $59/100000$. Add \eqref{eq:low-row} for the final assertion.
\end{proof}

This proof retains the selector until each positive norm is formed.
For the source's physical slots, its conservative pointwise envelopes
and moment inputs, with the same cumulative height allowance, apply
also to the required smooth derivative profiles. Under the abstract
alternative of stipulating the extra inputs, they must be assumed
uniformly for those derivative profiles as well.
A finer amplitude bin at one parameter is not automatically a bound
for those derivatives. The partial result does not cover the shorter
plain lengths near the simultaneous witness crossing.

\section{Relation to the exceptional-row count}\label{sec:application}

This section records a conditional application, including the source
of Hypothesis~\ref{hyp:count}. It is separate from the elementary
pair-count theorem.

\subsection{The existing scalar comparison}

Let $q$ be the physical-length-weighted mean prime amplitude, and set
$x=q/\delta$. In the relevant box $49/100\le x\le1/2$. Define
\begin{equation}\label{eq:scalar-params}
\begin{gathered}
\alpha=\frac56,\quad a_\delta=\alpha-\delta,\quad
B_x=2-\frac{8x}{9},\quad D_x=3-\frac{17x}{9},\\
P_x=B_x(1-x),\qquad J=a_\delta D_x+\delta P_x,\qquad
t_0=1+\frac{\delta P_x}{2J}.
\end{gathered}
\end{equation}
Here $D_x$ is a scalar and is distinct from the polynomial scale $D$.
The fixed-$\kappa=3/4$ comparison in the investigation gives
\begin{equation}\label{eq:Rstar}
R^*=1-\delta+\frac{a_\delta\delta P_x}{2J}
=1-\delta+a_\delta(t_0-1).
\end{equation}
Its input is the source's marginal moments and detector, using the
imported $\beta_*\le7/8$ statement to make this fixed $\kappa$ legal.
One must repeat that permitted fixed-parameter comparison rather than
quote the source's dynamic-$\kappa$ count with a negative
$\Delta=\beta_*-7/8$ allowance.

The count $\#\CC\ll U^{R^*+\ee}\HH^{A_0}$ concerns the full bin.
The detector may choose witness lengths separately to prove it.
There is therefore no circularity in using the count while asking for
a new mixed moment at prescribed $r,m$.

On $49/100\le x\le1/2$ one has $P_x/D_x\le2/5$. The function
$b/[2(a_\delta+b)]$ is increasing in $b>0$, so
\[
t_0-1
\le\frac{(2/5)\delta}{2(\alpha-\delta+(2/5)\delta)}
<\frac3{20}
\quad\left(\delta\le\frac{21}{50}\right).
\]
It follows that
\[
R^*\le1-\delta+\frac3{20}(\alpha-\delta)
=\frac98-\frac{23\delta}{20},
\]
which supplies Hypothesis~\ref{hyp:count}.

\subsection{Why both rebalances are required}

At scalar mean $q=\delta x$, the short inverse and plain count lines
are
\begin{align}
A_I(r)&=1-\delta[x+(1-x)r],\label{eq:AI}\\
S_t(r)&=1-\delta[4x/9+B_x(t-r)].\label{eq:St}
\end{align}
The long inverse count is
$L_\delta(t)=1-\delta+a_\delta(t-1)$. These expressions use the
source's original capacities, with independently small strict
decrements and whole-slot rounding.

Suppose the mixed estimate \eqref{eq:desired} is proved uniformly
where needed. On rows carrying the common detector spike,
\[
|M_rS_m|^2\gg U^{\delta(r+m)-o(1)}.
\]
Selecting prime slots of length close to $(1-r)/2$ with squared
gain at least $2qz-o(1)$ and applying the mixed moment improves
\eqref{eq:AI} by $\eta=1/5000$. It does not separately improve
\eqref{eq:St}. The inverse/plain crossing changes from
\[
r_*(t)=\frac{B_xt-5x/9}{D_x}
\quad\hbox{to}\quad
r_{\rm new}(t)=r_*(t)-\frac{\eta}{\delta D_x}.
\]
The optimized short count thus gains only $(B_x/D_x)\eta$.
Balancing this short count against the unchanged long count requires
a second adjustment:
\begin{equation}\label{eq:rebalance}
t_{\rm new}=t_0-\frac{B_x\eta}{J},\qquad
R_{\rm new}=R^*-\frac{a_\delta B_x}{J}\eta.
\end{equation}

\begin{proposition}[Conditional count payoff]\label{prop:payoff}
On the hard box, the gain factor satisfies
\begin{equation}\label{eq:factor}
\frac{a_\delta B_x}{J}\ge\frac{1091200}{2012413}.
\end{equation}
Consequently a uniform mixed saving $\eta=1/5000$ yields, before
the small application losses, count gain at least
\[
\frac{5456}{50310325}
=\frac1{10000}+\frac{169987}{20124130000}.
\]
\end{proposition}
\begin{proof}
Write $F=a_\delta B_x/J$. Direct differentiation gives
\[
\partial_\delta F
=-\frac{\alpha B_x^2(1-x)}{J^2}<0,\qquad
\partial_xF=\frac{a_\delta[(10/9)a_\delta+\delta B_x^2]}{J^2}>0.
\]
Its minimum is therefore at
$(\delta,x)=(21/50,49/100)$, where exact substitution gives
\eqref{eq:factor}. Multiply by $1/5000$.
\end{proof}

The requested length rectangle is sufficient for this conditional
application. Exact rational interval arithmetic gives
\[
1.106<t_{\rm new}<1.149,\qquad
0.701<r_{\rm new}<0.736.
\]
Within $0.70\le r\le0.74$, the detector relation
$m\ge t_{\rm new}-r-o(1)$ implies $m>0.366-o(1)$.
For $r\le0.70$ the old plain line is already below the improved
short envelope by at least $7/12500$; for
$0.74\le r<1$ the old inverse line has at least $13/25000$ slack,
including the cost $\eta$. Lengths $m\ge1/2$ use the original
unmarked plain moment, and $r\ge1$ uses the original long estimate.
Appendix~\ref{app:checks} identifies the exact record for these
intervals. A total count-level loss below $1/250000$ fits within
the surplus in Proposition~\ref{prop:payoff}.

\subsection{The smaller application domain}\label{sec:application-wedge}

The original rectangle is sufficient but larger than necessary for the
proposed count application. Put $m_{\rm new}=t_{\rm new}-r_{\rm new}$,
where $r_{\rm new}=r_*(t_{\rm new})-\eta/(\delta D_x)$, and retain
$R_{\rm new}$ from \eqref{eq:rebalance}. The old plain count at the
actual length is $P_{\rm plain}(m)=1-\delta(4x/9+B_xm)$. The exact
rebalance identities are
\begin{equation}\label{eq:wedge-crossing}
A_I(r_{\rm new})=R_{\rm new}+\eta,\qquad
P_{\rm plain}(m_{\rm new})=R_{\rm new}.
\end{equation}
Before the prescribed small losses, these old count bounds already
reach $R_{\rm new}$ if either
\[
r\ge r_{\rm new}+w,\quad w=\frac{\eta}{\delta(1-x)},
\qquad\hbox{or}\qquad m\ge m_{\rm new}.
\]
The detector support is $r+m\ge t_{\rm new}$ at this exponent
level. Thus the only region needing a new input has
\begin{equation}\label{eq:application-wedge}
a_r=r-r_{\rm new},\qquad b_m=m_{\rm new}-m,
\qquad 0<a_r<w,\qquad 0<b_m\le a_r.
\end{equation}
This statement localizes a sufficient application input; it does not
prove the full mixed moment outside this region by a row-count argument.

The exact width satisfies $1/1071\le w\le1/900$. On the whole
$(\delta,x)$ box, exact monotonicity gives
\begin{align*}
\frac{47749}{67660}&\le r_{\rm new}\le\frac{10261670}{14086891},\\
\frac{3470383}{8566666}&\le m_{\rm new}\le\frac{28646}{69475}.
\end{align*}
For example, with $F=a_\delta B_x/J$, the simplified identities
\[
r_{\rm new}=1-F/2-(1-F)w,\qquad
m_{\rm new}=\frac12-\frac{a_\delta}{J}
 \left(\frac{1-x}{2}-\frac\eta\delta\right)
\]
show that $r_{\rm new}$ increases in $\delta$ and decreases in $x$,
using the derivatives of $F$ above. The derivatives of $m_{\rm new}$
are positive on the box; the exact rational sign bounds are given in
the accompanying selector-energy note and checker. The leading wedge
lies in $0.7057<r<0.7296$ and $0.4039<m<0.4124$.
In particular it is disjoint from Proposition~\ref{prop:upper-band}'s
proved band.

A sufficient mixed saving inside this region can taper with $r$:
\begin{equation}\label{eq:tapered-target}
s(r)=\eta-\delta(1-x)(r-r_{\rm new})>0,\qquad
\sum_{\CC}|M_rS_mQ_I|^2
\ll_\ee U^{1+\delta m-s(r)+\ee}\HH^{A_0}.
\end{equation}
In the capacity limit, the simultaneous witness and prime lower bounds
then give count exponent $A_I(r)-s(r)=R_{\rm new}$. For an actual
prime product with squared gain exponent $2G_I$, the required saving
is instead $1-R_{\rm new}-\delta r-2G_I$, with all witness losses
included. A uniform $\eta$ saving is a stronger sufficient request
on the ideal wedge, before these losses.

One must use a buffered domain in the actual application. If $\ell_0$
bounds each marginal count loss and the detector support loss, then
the rows not already counted at exponent $R_{\rm new}$ satisfy
\begin{equation}\label{eq:buffered-wedge}
a_r<\frac{\eta+\ell_0}{\delta(1-x)},\qquad
b_m>-\frac{\ell_0}{\delta B_x},\qquad b_m\le a_r+\ell_0.
\end{equation}
The mixed theorem must be uniform on this outer region with its actual
capacity, profile and height losses. Those losses are chosen within
the existing conditional payoff surplus. No estimate for
\eqref{eq:tapered-target} on this region is proved here.

\subsection{Profile localization and the proposed boundary}

For completeness, the earlier profile refinement can be expressed
without changing the main theorem. Let $g^*$ be the decreasing
amplitude profile on total slot length $L_{\rm slot}$, with
\[
0\le g^*\le\delta/2,\qquad
q=L_{\rm slot}^{-1}\int_0^{L_{\rm slot}}g^*(w)\,dw,
\qquad L_{\rm slot}>1/5
\]
by a fixed supply margin. Put $G(v)=\int_0^v g^*(w)\,dw$, extended
constantly beyond $L_{\rm slot}$, and define
\begin{align*}
A_G(r)&=1-\delta r-2G((1-r)/2),\\
S_{G,t}(r)&=1-2\delta(t-r)-2G(2(1-2(t-r))/9),\\
F_G(t)&=\max_{t-1/2\le r\le1}\min\{A_G(r),S_{G,t}(r)\},\\
F_q(t)&=1-\delta+\frac{\delta P_x}{D_x}(3/2-t),\qquad
H_G(t)=F_q(t)-F_G(t).
\end{align*}
The preceding profile note \cite{records} shows, under the same
moment inputs and this whole-slot supply, that
$H_G(t_0)\ge1/5000$ gives full count gain
$31/260000>1/10000$ after retuning $t$, before arbitrarily small
capacity and whole-slot losses. Thus the new estimate can
be requested only on profiles with $H_G(t_0)<1/5000$.
Theorem~\ref{thm:main} applies to this smaller class, but does not
need that restriction.

The separate geometric calculation uses
\[
\ell=\frac16+\frac1{5000},\qquad b=\frac18,\qquad
l_x=\frac{5311}{15000},\quad l_y=\frac{3593}{7500},
\quad h=\frac{508}{625}.
\]
Its signal exponent is $C(s)=s-11/16$ and its low exponent is
$3749/20000$, giving the proposed matching boundary
$17499/20000=7/8-1/20000$. The saved continuous rational certificate
requires an additional row gain $1/10000$ in the remaining hard class.
Proposition~\ref{prop:payoff} explains why the proposed mixed saving
would supply that gain. This paragraph is a conditional implication
within the separately audited framework, not a zero-free theorem
proved here.

\section{Why the existing transformed proof does not close the gap}
\label{sec:obstruction}

The coefficient class matters even before the primitive conductor
restriction is imposed. On the core where the inverse ideal, the
plain ideal, and all physical primes are squarefree and mutually
coprime, write the total column as $v$. After separating its fixed
factor $\nu(v)$, the arithmetic coefficient has the form
$\mu_F(v)h(v)$, where
\begin{equation}\label{eq:h}
\begin{split}
h(v)=\sum_{kP\mid v}\mu_F(k)\mu_F(P)\,
A\!\left(\frac{\Nm v}{\Nm k\,\Nm P\,D}\right)B(\Nm k/N)\\
\hspace{3em}{}\times\prod_i a_i(p_i)W_i(\Nm p_i/P_i),
\qquad P=\prod_i p_i,
\end{split}
\end{equation}
with all original divisor, prime, and support restrictions retained.
This is an annular divisor sum with a free plain ideal. It is not
a product of a fixed number of individual prime-slot coefficients.

Under the substitutions of the source's inverse initialization,
the original squarefree column takes the form $v=tfn$.
The child coefficient therefore retains $h(tfn)$, which depends on
the canonical averaged label $f$ as well as the residual column $n$.
\cite[Lemma 17.2]{openai} expressly excludes an additional residual
coefficient, even one independent of the averaged labels.
A divisor bound on $h$ does not allow it to be inserted into that
signed theorem while preserving its conclusion.

One can make a more precise structural reduction without invoking that
theorem. Rename the transformed label $f$ as $g$, to distinguish it
from the ratio conductor. On the squarefree coprime support $v=tgn$,
assign each physical slot to its unique factor and write $P=P_0P_n$
according as it lies outside or inside $n$. Then the plain ideal splits
uniquely as
\[
k=k_0k_1,\qquad k_0\mid tg/P_0,\qquad k_1\mid n/P_n.
\]
The coefficient retains $\mu_F(k_0)\mu_F(k_1)$ and the coupled
profile $B(\Nm k_0\Nm k_1/N)$. Put $\lambda=\log_U\Nm k_1$.
After a common smooth separation, the source's coefficient identity
(17.72) permits factoring $n=k_1e$ and moving $k_1$ into the averaged
label $\widetilde g=gk_1$. Cauchy costs $U^{\lambda+\ee}$ in energy;
reindexing $(g,k_1)$ by $\widetilde g$ adds only divisor multiplicity.
This recovers the coefficient class for the separated positive child,
not a bound on its support.

In fact, let $L=r+m+z$, let $\beta$ be the norm exponent of $t$,
and let $V$ be that of $g$. The source's formal child centers are
\[
M=2L-1-2\beta,\qquad R=L-\beta-V,\qquad F=R+V.
\]
The change replaces $(R,V)$ by $(R-\lambda,V+\lambda)$, leaving
$F$ and the row center $M$ unchanged. If $z_n$ is the remaining
prime-slot length, its first canonical margin is still
\begin{equation}\label{eq:live-width}
F-M-\beta-z_n=1-L-z_n\le1-r-m\le-3/50.
\end{equation}
Even a hypothetical baseline canonical bound on that inadmissible
support would leave exponent $1+\lambda$ after this Cauchy step.
At $\lambda\approx m$ its excess over $K$ approaches
$(1-\delta)m+\eta\in[209/1000,1601/5000]$.
The detailed live-divisor note records the exact masks, slot assignments,
common separating profiles, and unequal-side normalization. This
complete-family calculation remains conditional on the selector gate
in Section~\ref{sec:selector-test}; it does not transform the sharp
kernel \eqref{eq:kernel}.

The widths fail as well. Appending the plain factor changes the
effective squarefree column length from $r+z$ to $r+m+z$ and adds
$2m$ to the first dual-row length in the source's equation (17.69).
Even if one ignores the forbidden coefficient, the formal canonical
margins at $z=(1-r)/2-\rho$, with a small strict decrement $\rho>0$,
become
\begin{equation}\label{eq:bad-widths}
2\rho-m,\qquad 2r-1-2m+8\rho.
\end{equation}
At $\rho=0$, these range in $[-1/2,-9/25]$ and
$[-3/5,-6/25]$. They are macroscopic deficits. Dropping every
prime slot would still require $r+m<1$, whereas the present rectangle
has $r+m\ge1.06$.

For a complete smooth row family, the source's masked Poisson
formula does give a signed nonzero-frequency expression for the
nonprincipal pairs. Bounding that enlarged positive norm is a
separate sufficient route. The selector in \eqref{eq:high-rows}
cannot simply be inserted into its Schwartz kernel or discarded
term by term from \eqref{eq:Tlarge}.

Finally, the equal-product centering used for two plain sums in
the source's Section 18 does not transfer just because an inverse
and a plain witness have common height. Their arithmetic coefficients
are different. Formal complete convolution $\mu_F*1$ also does not
annihilate a selected annular rectangle: for two appropriately
separated prime ideals, its semiprime coefficient can contain a
single nonzero M\"obius term. Any new recursion must propagate the
actual divisor coefficient or estimate its surviving signed
contribution.

\section{The remaining arithmetic question}

The strengthened task left by Theorem~\ref{thm:two-ratios} is
\eqref{eq:joint-open}, with the exact coefficient
\eqref{eq:tuple-coeff}, the row selector \eqref{eq:high-rows}, and
both large ratio-conductor restrictions in \eqref{eq:Tjoint}
retained. Proposition~\ref{prop:upper-band} proves a mixed subcase
for source-compatible physical slots. For the proposed count
application, the remaining new input can be confined to the buffered
region \eqref{eq:buffered-wedge}, with the actual saving budget
described there; that region is not covered by the upper-band result.
The reductions do not show that the surviving rows are sparse enough,
that their phases align, or that the signed kernel cancels.
In particular, exponent saturation of a normalized prime sum at
scale $P^{\delta/2}$ is far below its triangle-inequality scale
$P^{1/2+o(1)}$ on the present box; it supplies no elementary
near-equality argument.

A proof of \eqref{eq:joint-open} would complete the proposed localized
mixed input and permit the two rebalances of
Section~\ref{sec:application}. At present it remains an explicit
open estimate.

\appendix
\section{Dependencies and reproducible arithmetic}\label{app:checks}

The results in this draft have three distinct statuses.
\begin{enumerate}
\item Lemmas~\ref{lem:ratio}--\ref{lem:pair-count} and
Proposition~\ref{prop:small} are elementary coefficient, local-character,
and ideal-counting statements proved here.
\item Theorems~\ref{thm:main} and \ref{thm:two-ratios} are conditional on
Hypotheses~\ref{hyp:inverse}--\ref{hyp:count}. Their application uses
the source's buffered, marked-moment, detector, and fixed-parameter
row-count framework. Proposition~\ref{prop:upper-band} additionally
uses the source-compatible fourth moment and buffered pointwise inputs
stated in Section~\ref{sec:mixed-subcase}. Those deep results are not
re-proved here.
\item The signed estimates \eqref{eq:open} and \eqref{eq:joint-open},
the full mixed bound
\eqref{eq:desired}, and the proposed improved boundary remain unproved.
\end{enumerate}

The investigation directory accompanying this manuscript retains
small, standard-library exact-arithmetic checks:
\begin{itemize}
\item \path{numerics/check_mixed_kernel.py} checks the local
valuation/mask identities on 6,561 finite valuation pairs and the
rational removal margins and width deficits. The infinite pair-count
bound is proved in Section~\ref{sec:pairs}, not inferred from this
enumeration.
\item \path{numerics/check_joint_witness_payoff.py} checks symbolic
derivative identities, the two rebalances, the uniform length
intervals, and their rational slack.
\item \path{numerics/check_selector_feasibility.py} checks the
principal-row inverse budget, the mean-zero subcase margin, and
the unchanged widths after the live plain divisor enters an averaged
label. It does not establish nonvanishing profile masses or bound
the remaining mixed correlation.
\item \path{numerics/check_selector_preserving.py} checks both
adaptive conductor cutoffs, their local zero masks, and the upper
plain-band loss budget. \path{numerics/check_selector_energy.py}
checks the exact rebalance identities, continuous monotonicity bounds,
and the smaller sufficient application domain. Neither checker proves
the correlation estimate within that domain.
\item \path{numerics/check_localized_target.py} certifies the
conditional geometric payoff by positive rational Bernstein
coefficients on a continuous rectangle cover. Its hard-box test
assumes the additional row gain; it does not prove that gain.
\end{itemize}
Each script prints its small retained JSON record. Decimal displays
are explanatory; signs and margins are decided by exact arithmetic.
No finite check here estimates the remaining arithmetic character
correlation.

The consulted source PDF is identified by SHA-256
\begin{center}\small
\path{8fe93046f8cf5ef1ba5969c89addc02d76311adc4ee907509ff9cd96f7ec99e7}.
\end{center}
It is cited rather than copied into the repository. Same-model
parallel reviews checked the displayed reductions and their
application hypotheses. They are not independent specialist
refereeing or formal verification.

\section*{Acknowledgements and authorship status}
This working manuscript was drafted for Edward Baker with substantial
assistance from GPT-6 (Codex), including exposition, algebraic
derivations, exact-arithmetic checks, and separate agent reviews.
The exact serving variant and configured reasoning effort were not
exposed and are not inferred. The draft has not been represented as
author-approved, peer-reviewed, or formally verified. Responsibility
for any eventual submission and its mathematical claims remains with
the named human author after review.

\begin{thebibliography}{9}
\bibitem{openai}
OpenAI,
\emph{The Quasi-Riemann Hypothesis: A Zero-Free Half-Plane
$\Rea s>7/8$},
preprint, 30 September 2026.
\url{https://github.com/openai/math/blob/main/preprints/The-Quasi-Riemann-Hypothesis-September-30-2026/paper.pdf}.

\bibitem{records}
Research records prepared for Edward Baker with GPT-6 (Codex)
assistance, \emph{Quasi-RH character amplification},
8 October 2026. Accompanying investigation notes:
\emph{The first mixed-moment calculation: one remaining signed sum};
\emph{Mixed-witness kernel: small conductors are already harmless};
\emph{Primitive-conductor localization of the remaining mixed moment};
and \emph{Localized joint-witness target and its conditional payoff}.
The application also uses the amplitude-profile, joint-witness, and
geometry records at
\path{notes/AMPLITUDE_PROFILE_REDUCTION_20261008.md},
\path{notes/JOINT_WITNESS_REDUCTION_20261008.md}, and
\path{notes/GEOMETRY_OPTIMIZATION_20261008.md}.
See the investigation's \path{README.md} and
\path{numerics/README.md} for the complete index and exact checks.
\end{thebibliography}
\end{document}
